\chapter{Stokes Structures}
\begin{comment}
  Matrizen Sicht:
  \begin{itemize}
    \item \cite[Chapter 3]{boalch}
      \\ largely following:
      \begin{itemize}
        \item[\textbf{8}] D.G. Babbitt and V.S. Varadarajan.
        \textbf{Formal reduction theory of meromorphic differential equations:
          a group theoretic view.}
        \texttt{euclid.pjm.1102720203.pdf}
        \item[\textbf{11}] W. Balser, W.B. Jurkat, and D.A. Lutz.
        \textbf{Birkhoff invariants and Stokes’ multipliers for meromorphic
          linear differential equations.}
        \item[\textbf{40}] M. Jimbo, T. Miwa, and Kimio Ueno.
        \textbf{Monodromy preserving deformations of linear differential
          equations with rational coefficients I.}
        \item[\textbf{43}] M. Loday-Richaud
        \textbf{Stokes phenomenon, multisummability and differential Galois
          groups.}
        \item[\textbf{50}] J. Martinet and J.P. Ramis.
        \textbf{Elementary acceleration and multisummability.}
      \end{itemize}
    \item \cite{thboalch}
    \item \cite{van2003galois}
    \item Marius van der Put, Kyoshi Saito => Diff.Galois Theory
  \end{itemize}
  Moderne Sicht (sheaf-view):
  \begin{itemize}
    \item Malgrange
    \item \cite{sabbah_cimpa90} and
    \item \cite{sabbah2007isomonodromic}
    \item \cite{sabbah2009introduction} and \cite{sabbah2013introduction}
  \end{itemize}
\end{comment}

\section{Matrix view: \cite{boalch} and \cite{thboalch}} %{{{
Let
\begin{itemize}
  \item $z$ be a coordinate on $\D$ vanishing at $0$ and
  \item $d-A^0$ be a diagonal meromorphic connection
    \begin{itemize}
      \item on the trivial rank $n$ vector bundle
        \begin{itemize}
          \item over the unit disc $\D\subset\C$
        \end{itemize}
      \item with a pole of order $k\geq2$ at $0$
        \begin{itemize}
          \item and no other poles.
        \end{itemize}
    \end{itemize}
    thus $A^0=dQ+\Lambda^0\frac{dz}{z}$ where
    \begin{itemize}
      \item $\Lambda^0$ is a constant diagonal matrix and
      \item $Q$ is a diagonal matrix of meromorphic functions.
        \begin{itemize}
          \item write $Q=\diag(g_1,\dots,q_n)$
        \end{itemize}
    \end{itemize}
\end{itemize}
Define $q_{ij}(z)$ to be the leading term of $q_i-q_j$.
\begin{itemize}
  \item Thus if $q_i-q_j=a/z^{k-1}+b/z^{k-2}+\dots$ then $q_{ij}=a/z^{k-1}$.
\end{itemize}
\subsection{Basic definitions}
\begin{defn}[Definition 3.2]
The \emph{anti-Stokes} directions $\A\subset S^1$ are the directions $d\in
S^1$ such that for some $i \neq j$: $q_{ij}(z)\in\R_{<0}$ for $z$ on the ray
specified by $d$.
\end{defn}
\begin{center}
  \begin{tikzpicture}[scale=2]
    \node[label=below left:$0$] (zero) at (0,0) {};
    \draw[blue] (zero) circle (1cm);
    \draw[red!60!black,thick,path fading=east] (0,0) -- +({cos( 33 )*2},{sin( 33 )*2});
    \fill[blue!20!white] ({cos( 33 )},{sin( 33 )}) circle (1pt);

    \node[blue] at (-0.9,0.6) {$S^1$};
    \node at (-1.3,-0.7) {$\C$};
    \node[blue!60!black,right] at ({cos( 33 )},{sin( 33 )}) {$d$};
    \node[red!60!black] at (2.2,.9) {ray specified by $d$};
    \fill (zero) circle (1pt);
  \end{tikzpicture}
\end{center}
\begin{rem}
  \begin{itemize}
    \item We have $\frac{\pi}{k-1}$ rotational symmetry
      \begin{itemize}
        \item if $q_{ij}(z)\in\R_{<0}$ then
          $q_{ij}(z\exp(\frac{\pi\sqrt{-1}}{k-1}))\in\R_{<0}$
      \end{itemize}
    \item If $q_{ij}(z)\in\R_{<0}$ then $q_{ji}(z)\in\R_{>0}$
      \begin{itemize}
        \item in any arc $U \subset S^1$ subtending angle $\frac{\pi}{k-1}$
          there are at most $\frac{n(n-1)}{2}$ anti-Stokes directions.
      \end{itemize}
  \end{itemize}
\end{rem}
\begin{defn}
  \begin{itemize}
    \item Define the \emph{set of all anti stokes directions}
      \begin{tikzpicture}[scale=3]
        \fill[blue!20!white] (0,0) circle (1pt);
      \end{tikzpicture}
      $\in\A\subset S^1$,
      \begin{itemize}
        \item
          \textbf{Labeling convention:}
          Choose a point\footnote{In~\cite{thboalch} a sector is chosen, which
          does not contain any anti Stokes direction}
          \textcolor{yellow!60!black}{$p$}$\in \D$.
          Label the first anti-Stokes ray when turning in a positive sense from
          $p$ as $d_1$ and label the subsequent rays $d_2,\dots,d_r$ in turn.
          \begin{rem}
            The choice of labeling together with the choice of $\log(z)$ is
            given by a lift $\tilde p$ of $p$ to the universal cover of $D^*$.
          \end{rem}
      \end{itemize}
  \end{itemize}
  \begin{itemize}
    \item
    \item $r:=\#\A$
      \begin{itemize}
        \item $r$ is divisible by $2(k-1)$
          \begin{itemize}
            \item define $l:=\frac{r}{2k-2}$
          \end{itemize}
      \end{itemize}
    \item $\textbf{d}:=(d_i,d_{i+1},d_{i+2},\dots,d_{i+l})\subset\A$ is a
      \emph{half-period}\footnote{A $l$-tuple of \textbf{consecutive}
      anti-Stokes directions}
    \item Define the open sector $\Sect(d_i,d_j)$
      \begin{itemize}
        \item The radius \textcolor{green!40!black}{$r$} will be taken
          sufficiently small when required later.
        \item Write $\Sect_i:= \Sect(d_i,d_{i+1})$ the ‘ith sector’.
          \begin{itemize}
            \item Note \textcolor{yellow!60!black}{$p$} $\in\Sect_r= \Sect_0$
          \end{itemize}
      \end{itemize}
    \item $\widehat{\Sect}_i:=
      \Sect(d_i-\frac{\pi}{2k-2},d_{1+1}+\frac{\pi}{2k-2})$ the ‘ith
      supersector’
      \begin{itemize}
        \item The rays bounding these supersectors are usually refered as
          \textcolor{red!40!black}{‘Stokes rays’}.
      \end{itemize}
  \end{itemize}
  \begin{center}
    \begin{tikzpicture}[scale=3]
      \node[] (zero) at (0,0) {};
      \fill[fill=green!20!white] (0,0) -- (1,0) arc (0:60:1.0cm) -- cycle;
      \draw[blue] (zero) circle (1cm);

      \foreach \w/\str in {10/$d_1\in S^1$,
                           20/$d_2$,
                           45/$d_3$,
                           55/$d_l$}
      {\draw[thick,purple!\w!blue,path fading=west]
          (0,0) -- +({cos( \w )},{sin( \w )}) node[right] {\str};
       \fill[blue!20!white] ({cos( \w )},{sin( \w )}) circle (1pt);
       \foreach \sep in {60,120,180,240,300}
       {\draw[green!20!white,thick] (zero) -- +({cos( \sep )},{sin( \sep )});
        \draw[purple!\w!blue] (0,0) -- +({cos( \w + \sep )},{sin( \w + \sep )});
        \fill[blue!20!white] ({cos( \w + \sep )},{sin( \w + \sep )}) circle (1pt);
       }
      };

      \foreach \sep/\str in {0/$1$
                            ,60/$2$
                            ,120/$k-1$
                            ,180/$4$
                            ,240/$5$
                            ,300/$2(k-1)$}
      {\node[green!40!black]
        at ({.6 * cos( \sep + 30 )},{.5 * sin( \sep + 30)}) {\str};
      };

      \fill[yellow!60!black] (0.8,0.07) circle (1pt);
      \node[yellow!60!black,right] at (0.8,0.07) {$p$};

      \fill (zero) circle (1pt);
    \end{tikzpicture}
  \end{center}
  \begin{center}
    \begin{tikzpicture}[scale=3]
      \node[] (zero) at (0,0) {};
      \draw[blue] (zero) circle (1cm);
      \fill[fill=green!20!white]
        (0,0) -- +({cos(10) * 0.7},{sin(10) * 0.7}) arc (10:45:0.7cm) -- cycle;
      \draw[->,thick,green!40!black] ({cos(10) * 0.7},{sin(10) * 0.7}) arc (10:45:0.7cm);

      \foreach \w/\str in {10/$d_1\in S^1$,
                           20/$d_2$,
                           45/$d_3$,
                           55/$d_l$}
      {\draw[thick,path fading=west] (0,0) -- +({cos( \w )},{sin( \w )}) node[right] {\str};
       \foreach \sep in {60,120,180,240,300}
       {\draw[gray] (0,0) -- +({cos( \w + \sep )},{sin( \w + \sep )});}
      };

      \fill (zero) circle (1pt);

      \fill[yellow!60!black] (0.8,0.07) circle (1pt);
      \node[yellow!60!black,right] at (0.8,0.07) {$p$};

      \node[green!40!black] at (1.2,0.6) {$\Sect(d_1,d_3)$};
      \node[green!40!black] at (0.4,0.02) {$r$};

    \end{tikzpicture}
  \end{center}
  \begin{center}
    \begin{tikzpicture}[scale=3]
      \node[] (zero) at (0,0) {};
      \draw[blue] (zero) circle (1cm);

      \fill[fill=green!20!white] (0,0) -- ({cos( 15 )},{sin( 15 )}) arc
        (15:85:1) -- cycle;

      \fill[fill=red!60!black] (0,0) -- ({cos( 15 )*0.5},{sin( 15 )*0.5}) arc
        (15:45:0.5) -- cycle;
      \fill[fill=red!60!black] (0,0) -- ({cos( 55 )*0.5},{sin( 55 )*0.5}) arc
        (55:85:0.5) -- cycle;

      \fill[fill=green!20!white] (0,0) -- ({cos( 15 )*0.4},{sin( 15 )*0.4}) arc
        (15:85:0.4) -- cycle;

      \node[green!40!black] at (1,1) {$\widehat\Sect_i$};

      \node[] (lft) at ({cos( 85 )},{sin( 85 )}) {};
      \node[] (rgt) at ({cos( 15 )},{sin( 15 )}) {};

      \draw[->,red!40!black] ({cos( 45 )},{sin( 45 )})
        to [out=35, in=25] (rgt);

      \draw[->,red!40!black] ({cos( 55 )},{sin( 55 )})
        to [out=65, in=75] (lft);

      \draw[thick,red!40!black,path fading=west] (0,0) -- +({cos( 85 )},{sin( 85 )});
      \draw[thick,red!40!black,path fading=west] (0,0) -- +({cos( 15 )},{sin( 15 )});
      \fill[red!40!black] ({cos( 85 )},{sin( 85 )}) circle (1pt);
      \fill[red!40!black] ({cos( 15 )},{sin( 15 )}) circle (1pt);

      \foreach \w/\str in {10/$d_1$,
                           20/$d_2$,
                           45/$d_i$,
                         55/$d_{i+1}$}
      {\draw (0,0) -- +({cos( \w )},{sin( \w )}) node[right] {\str};
       \fill[blue!20!white] ({cos( \w )},{sin( \w )}) circle (.7pt);
       \foreach \sep in {60,120,180,240,300}
       {\draw (0,0) -- +({cos( \w + \sep )},{sin( \w + \sep )});
        \fill[blue!20!white]
          ({cos( \w + \sep )},{sin( \w + \sep )}) circle (.7pt);
       }
      };

      \fill[yellow!60!black] (0.8,0.07) circle (1pt);
      \node[yellow!60!black,right] at (0.8,0.07) {$p$};

      \fill (zero) circle (1pt);
    \end{tikzpicture}
  \end{center}
\end{defn}
\begin{defn}
  \begin{itemize}
    \item $\Roots(d):=\{(ij)\mid g_{ij}(z)\in \R_{<0} \text{ along } d \}$
    \item The \emph{multiplicity} $\Mult(d)$ of $d$ is the number of roots
      supporting $d$ ($\#\Roots(d)$).
      \begin{rem}
        $n(n-1)/2=\Mult(d_1)+\cdots+\Mult(d_l)$
      \end{rem}
    \item The group of \emph{Stokes factors} associated to $d$ is the group
    \[
      \SSto_d(A^0) := \{K \in G \mid (K)_{ij}
        =\delta_{ij} \text{ unless } (ij) \text{ is a root of } d\}.
    \]
    \begin{itemize}
      \item $i=j$ \Rightarrow $(ij)$ is not a root of $d$. There are $1$nes
        on the diagonal.
      \item is a unipotent subgroup of $G=GL_n(\C)$
        \begin{itemize}
          \item of dimension equal to $\Mult(d)$
        \end{itemize}
    \end{itemize}
  \end{itemize}
\end{defn}
\begin{lem}[3.2]
  \TODO[move down?]
  let $\textbf{d}$ be a half-period.
  \begin{enumerate}
    \item \textbf{The product of the corresponding groups of Stokes factors} is
      isomorphic
      \begin{itemize}
        \item as a variety,
        \item via the product map,
      \end{itemize}
      to the \textbf{$\bf P$ conjugate of $\bf U_+$}\footnote{Is a subgroup of
      $G$.}:
      \[
        \prod_{d\in\textbf{d}}\SSto_d(A^0)\cong PU_+P^{-1} ;
        \qquad
        (K_1,\dots,K_l)\mapsto K_l\dots K_2K_1\in G
      \]
      where
      \begin{itemize}
        \item $P\in G$ the permutation Matrix associated to $\textbf{d}$ given
          by $(P)_{ij}=\delta_{\pi(i)j}$ where
          \begin{itemize}
            \item $\pi$ is the permutation of $\{1,\dots,n\}$ corresponding to
            \[
              \pi(i)<\pi(j)
              \qquad \Leftrightarrow \qquad
              q_i\underset{\textbf{d}}{<}q_j
            \]
            where
            \begin{itemize}
              \item the total ordering is defined as
                \[
                  q_i\underset{\textbf{d}}{<}q_j \qquad :\Leftrightarrow
                    \qquad (ij) \text{ is a root of some } d\in\textbf{d}
                \]
            \end{itemize}
          \end{itemize}
      \end{itemize}
    \item The map
      \[
        \prod_{d\in\A}\SSto_d(A^0)\cong (U_+\times U_-)^{k-1} ;
        \qquad
        (K_1,\dots,K_r)\mapsto (S_1,\dots,K_{2k-2})
      \]
      where
      \begin{itemize}
        \item $S_i:=P^{-1}K_{il}\dots K_{(i-1)l+1}P\in U_{+/-}$ if $i$ is
        odd / even
      \end{itemize}
      from the \textbf{product of all the groups of
      Stokes factors}, is
      \begin{itemize}
        \item an isomorphism of varieties.
      \end{itemize}
  \end{enumerate}
\end{lem}

\subsection{Main theoreme: for Stokes factors}
\begin{thm}[1.22 in \cite{thboalch}]
  \begin{itemize}
    \item There is a natural isomorphism
      \[
        \cH(A^0)\cong\prod_{d\in\A}\Sto_d(A^0)
      \]
      and
    \item for each choice of $\log(z)$ in the direction $d$
      \begin{itemize}
        \item the Stokes group $\Sto_d(A^0)$ has
          \begin{itemize}
            \item a faithfull representation $\rho$ on $\C^n$
          \end{itemize}
          inducing an isomorphism
          \[
            \rho:\Sto_d(A^0)\to\SSto_d(A^0).
          \]
      \end{itemize}
  \end{itemize}
  In particular each $\Sto_d(A^0)$ and therefore $\cH(A^0)$ is a unipotent Lie
  group and the complex dimension of $\cH(A^0)$ is $(k-1)n(n-1)$ where
  \begin{itemize}
    \item $k$ is the order of the pole of $A^0$ and
    \item $n$ is the rank.
  \end{itemize}
\end{thm}
\begin{thm}[3.1, or Prop 1.24 in \cite{thboalch}]
  Let
  \begin{itemize}
    \item $\hat{F}\in G\llbracket z\rrbracket$
      \begin{itemize}
        \item such that $A:=F[A_0]$ has convergent entries.
      \end{itemize}
  \end{itemize}
  Set the radius of the sectors $\Sect_i$, $\widehat\Sect_i$ to be less than
  the radius of convergence of $A$.
  Then:
  \begin{enumerate}
    \item On each sector $\Sect_i$ there is a \textbf{canonical} way to choose
      \begin{itemize}
        \item $\Sigma_i(\hat F)$
          \begin{itemize}
            \item an invertible $n\times n$ matrix of holomorphic functions
          \end{itemize}
      \end{itemize}
      such that $\Sigma_i(\hat F)[A^0]=A$.
    \item $\Sigma_i(\hat F)$ can be analytically continued to the supersector
      $\widehat\Sect_i$ and then
      \[
        \Sigma_i(\hat F)
        \qquad\text{is \textbf{asymptotic} to}\qquad
        \hat F
      \]
      \begin{itemize}
        \item at $0$
        \item within $\widehat\Sect_i$.
      \end{itemize}
    \item If $g\in G\{z\}$ and $t\in T$ then
      $\Sigma_i(g\circ\hat F \circ t^{-1})=g\circ\Sigma_i(\hat F)\circ t^{-1}$.
      \begin{prop}[Prop 1.24 in \cite{thboalch}]
        If $\hat H\in\hat G$ is the taylor series at $0$ of an analytic gauge
        transformation $H\in G\{z\}$ then
        \[
          \Sigma_i(\hat H \hat F)=H\Sigma_i(\hat F)
          \qquad
          \Sigma_i(\hat F \hat H)=\Sigma_i(\hat F)H
        \]
      \end{prop}
  \end{enumerate}
  \begin{rem}
    \begin{itemize}
      \item The point is that on a narrow sector there are generally many
        holomorphic isomorphisms between $A_0$ and $A$ which are asymptotic to
        $\hat F$ and one is being chosen in a canonical way
      \item $\Sigma_i(\hat F)$ is in fact uniquely characterised by property 2.
    \end{itemize}
  \end{rem}
\end{thm}
\paragraph{Now define the map $\cH(A^0)\to\prod_{d\in\A}\Sto_d(A^0)$:} choose
\begin{itemize}
  \item $\hat F\in\hat G(A^0)$
    \begin{itemize}
      \item representing an element
        $\cH(A^0)\cong G\{z\}\backslash\hat G(A^0)$
    \end{itemize}
    and
  \item an anti-Stokes direction $d\in\A$.
\end{itemize}
The sums of $\hat F$ on the two sectors adjacent to $d$ may be analytically
continued across $d$, and they will generally be different on the overlap. Thus
to each anti-Stokes direction $d=d_i$ there is
\begin{itemize}
  \item an associated automorphism
    \[
      \kappa_{d_i}:=(\Sigma_i(\hat F))^{-1}\Sigma_{i-1}(\hat F)
    \]
    \begin{itemize}
      \item describing how the sums of $\hat F$ differ on both sides of $d_i$
      \item it is a solution of $\Hom(A^0,A^0)$ asymptotic to $1$ on a
        sectorial neighbourhood of $d_i$
      \item The third part of the Proposition implies that the automorphism
        $\kappa_d$ only depends on the $G\{z\}$ orbit of $\hat F$.
    \end{itemize}
\end{itemize}
\begin{defn}[\cite{thboalch} 1.25]
  The \emph{Stokes group} $\Sto_d(A^0)$ of such automorphisms that arise as we
  vary the choice of $\hat F$:
  \[
    \Sto_d(A^0):=
      \left\{\kappa_d=(\Sigma_i(\hat F))^{-1}\Sigma_{i-1}(\hat F)
        \mid \hat F\in\hat G(A^0)\right\}
  \]
\end{defn}
Taking all such automorphisms gives a map
\[
  \hat G(A^0)\to\prod_{d\in\A}\Sto_d(A^0);
    \hat F\mapsto(\kappa_{d_1},\dots,\kappa_{d_r})
\]
inducing a well defined map
\[
  \cH(A^0)\to\prod_{d\in\A}\Sto_d(A^0) \,.
\]
\begin{comment}
  Using the Proposition:
  \begin{align*}
    \kappa_{d_i}&=(\Sigma_i(\hat F))^{-1}\Sigma_{i-1}(\hat F)
    \\&=(\Sigma_i(\hat F))^{-1}g^{-1}g\Sigma_{i-1}(\hat F)
    \\&=(g\Sigma_i(\hat F))^{-1}g\Sigma_{i-1}(\hat F)
    \\&=(\Sigma_i(g\hat F))^{-1}\Sigma_{i-1}(g\hat F)
  \end{align*}
\end{comment}
\paragraph{The faithfull representation of $\Sto_d(A^0)$} arises since the
solutions of $\Hom(A^0,A^0)$ are known explicitly:
\\Let
\begin{itemize}
  \item $d$ be a direction and
  \item a choice of $\log(z)$ in the direction $d$ be given.
\end{itemize}
We get a genuine fundamental solution $z^\Lambda e^Q$ of $A^0$ there.
\begin{paracol}{2} %%%%%%%%%%%%%%%%%%%%%%%%%%%%%%%%%%%%%%%%%%%%%%%%%%%%%%%%%%%%
  The map
  \[
    \rho:\kappa\to K:=e^{-Q}z^{-\Lambda}\kappa z^\Lambda e^Q
  \]
  relates
  \begin{itemize}
    \item solutions $\kappa$ of $\Hom(A^0,A^0)$
  \end{itemize}
  to
  \begin{itemize}
    \item a konstant matrix $K\in\End(E)$
      \begin{itemize}
        \item that is, to solutions of $\Hom(0,0)$.
      \end{itemize}
  \end{itemize}
  \begin{lem}
    The Stokes factor $K_i$ is in the group $\SSto_{d_i}(A^0)$.
  \end{lem}
  \switchcolumn %%%%%%%%%%%%%%%%%%%%%%%%%%%%%%%%%%%%%%%%%%%%%%%%%%%%%%%%%%%%%%%
  \begin{defn}[\cite{boalch}: def 3.5]
    Fix data $(A^0,z,\tilde\rho)$\footnote{as above.}.
    The
    \begin{itemize}
      \item \emph{Stokes factors}
        \begin{itemize}
          \item of a compatibly framed connection
            $(A,g^0)\in\widehat\Syst(A^0)$
        \end{itemize}
    \end{itemize}
    are
    \[
      K_i:=e^{-Q}z^{-\Lambda^0}\kappa z^{\Lambda^0} e^Q
    \]
    using
    \begin{itemize}
      \item the choice of $\log(z)$ along $d_i$, where
        $\kappa_i:=\Sigma_i(\hat F)^{-1}\circ\Sigma_{i-1}(\hat F)$.
    \end{itemize}
  \end{defn}
  Since
  \begin{itemize}
    \item $z^{\Lambda^0}e^Q$ is a fundamental solution of $A^0$
  \end{itemize}
  we have
  \begin{itemize}
    \item $d(K_i)=0$
      \begin{itemize}
        \item The Stokes factors are constant invertible matrices.
      \end{itemize}
  \end{itemize}
  A usefull equivalent definition is:\dots
\end{paracol} %%%%%%%%%%%%%%%%%%%%%%%%%%%%%%%%%%%%%%%%%%%%%%%%%%%%%%%%%%%%%%%%%
\begin{defn}[1.27]
  Given $\tilde p$\footnote{a choice of labeling and $\log(z)$}.
  The \emph{Stokes factors} 
  \begin{itemize}
    \item of a compatibly framed system $(A,g)$
    \item with associated normal Form $A^0$
  \end{itemize}
  are
  \[
    \textbf{K}:=(K_1,\dots,K_r)\in\prod_{d\in\A}\SSto_d
  \]
  \begin{itemize}
    \item $r$-tuple of matrices
    \item $K_i=\rho(\kappa_i)=e^{-Q}z^{-\Lambda}\kappa z^\Lambda e^Q$
      \begin{itemize}
        \item using the choice of $\log(z)$ along $d$.
      \end{itemize}
  \end{itemize}
\end{defn}
\subsubsection{A usefull way of thinking about these}
Fix
\begin{itemize}
  \item a choice of labelling and
  \item a choice of $\log(z)$.
\end{itemize}
\begin{defn}
  If
  \begin{itemize}
    \item $(A,g)$ is a compatibly framed system
      \begin{itemize}
        \item with associated formal normal form $A^0$
      \end{itemize}
  \end{itemize}
  then the \emph{cononical fundamental solution of $A$} on $\Sect_i$ is
  \[
    \Phi_i:=\Sigma_i(\hat F)z^\Lambda e^Q
  \]
  where
  \begin{itemize}
    \item $z^\Lambda$ uses the given choice of $\log(z)$ on $\Sect_i$.
  \end{itemize}
  \begin{rem}[???]
    \begin{itemize}
      \item $\Phi_i$ is a solution of $\Hom(0,A)$
      \item $\Phi_i$ is asymptotic to $\hat Fz^\Lambda e^Q$
        \begin{itemize}
          \item at $0$
          \item within $\widehat \Sect_i$
        \end{itemize}
        when continued without any winding.
    \end{itemize}
  \end{rem}
\end{defn}
Then we get:
\begin{paracol}{2} %%%%%%%%%%%%%%%%%%%%%%%%%%%%%%%%%%%%%%%%%%%%%%%%%%%%%%%%%%%%
  \begin{lem}
    If $\Phi_i$ is continued across the anti-Stokes ray $d_{i+1}$ then on
    $\Sect_{i+1}$:
    \[
      \Phi_i=\Phi_{i+1}K_{i+1}
      \qquad \text{ for } i=1,\dots,r-1 \text{ and}
    \]
    \[
      \Phi_r=\Phi_1K_1M_0
    \]
    where
    \begin{itemize}
      \item $M_0=e^{2\pi\sqrt{-1}\Lambda}$ is the formal monodromy.
    \end{itemize}
  \end{lem}
  \switchcolumn %%%%%%%%%%%%%%%%%%%%%%%%%%%%%%%%%%%%%%%%%%%%%%%%%%%%%%%%%%%%%%%
  \begin{defn}[\cite{boalch}: def 3.6]
    If $\Phi_i$ is continued across the anti-Stokes ray $d_{i+1}$ then on
    $\Sect_{i+1}$ we have:
    \[
      K_{i+1}:=\Phi_{i+1}^{-1}\circ\Phi_i
    \]
    for all $i$ except $K_1:=\Phi_{r+1}^{-1}\circ\Phi_r\circ M_0^{-1}$
    where
    \begin{itemize}
      \item $M_0=e^{2\pi\sqrt{-1}\Lambda}$ is the so-called `formal monodromy'.
    \end{itemize}
  \end{defn}
  \begin{lem}
    The Stokes factor $K_i$ is in the group $\SSto_{d_i}(A^0)$.
  \end{lem}
\end{paracol} %%%%%%%%%%%%%%%%%%%%%%%%%%%%%%%%%%%%%%%%%%%%%%%%%%%%%%%%%%%%%%%%%
Thus, in summary, a compatibly framed system $(A,g)$ has canonical
fundamental solutions $\Phi_i$ and the Stokes factors express how these
differ and moreover they encode the moduli of $(A,g)$.

\subsection{Main theoreme: for Stokes matrices}
\begin{paracol}{2} %%%%%%%%%%%%%%%%%%%%%%%%%%%%%%%%%%%%%%%%%%%%%%%%%%%%%%%%%%%%
  \switchcolumn %%%%%%%%%%%%%%%%%%%%%%%%%%%%%%%%%%%%%%%%%%%%%%%%%%%%%%%%%%%%%%%
  Thus we can define the Stokes matrices of $(A,g_0)\in\hat\Syst(A^0)$
  \[
    S_i:=P^{-1}K_{il}\dots K_{(i-1)l+1}P\in U_{+/-}
    \qquad
    \text{if } i \text{ is odd / even}
  \]
  where
  \begin{itemize}
    \item $i=1,\dots,2k-2$ and
    \item $P$ is the permutation matrix associated to the half-period
      $(d_1,\dots,d_l)$.
  \end{itemize}
\end{paracol} %%%%%%%%%%%%%%%%%%%%%%%%%%%%%%%%%%%%%%%%%%%%%%%%%%%%%%%%%%%%%%%%%
\begin{prop}[\cite{thboalch}: 1.35]
  $P$ is the (unique) permutation matrix, such that the multiplication map
  \begin{align*}
    \SSto_{il}(A^0)\times\dots\times\SSto_{(i-1)l+1}(A^0)
    &\to \GL_n(\C)
  \\(K_{il},\dots,K_{(i-1)l+1})&\mapsto P^{-1}K_{il}\cdots K_{(i-1)l+1}P
  \end{align*}
  is
  \begin{itemize}
    \item a diffeomorphism onto
      $
      \begin{cases}
        U_+ & \text{if $i$ is odd}
      \\U_- & \text{if $i$ is even}
      \end{cases}
      $
  \end{itemize}
  where
  \begin{itemize}
    \item $i=1,\dots$ and
    \item all depends on the ordering $d_1,\dots,d_r$ of the anti-stokes directions.
  \end{itemize}
  \begin{defn}
    The \emph{$i$th Stokes matrix}
    \begin{itemize}
      \item with respect to a fixed labelling and $\log(z)$ choice
    \end{itemize}
    is the following product of $l$ Stokes factors
    \[
      S_i:=K_{il}\cdots K_{(i-1)l+1}\in PU_{\pm}P{-1}\subset\GL_n(\C)
    \]
  \end{defn}
\end{prop}
\begin{thm}[\cite{boalch}: Thm 3.2]
  Fix the data $(A^0,z,\tilde p)$ as above.
  Then 
  \begin{itemize}
    \item the \emph{`local monodromy map'}
      \begin{itemize}
        \item taking the stokes matrices
      \end{itemize}
  \end{itemize}
  induces a bijection
  \begin{align*}
    \cH(A^0)&\overset{\cong}{\to}(U_+\times U_-)^{k-1}
  \\ \left[(A,g_0)\right] &\mapsto(S_1,\dots,S_{2k-2})
  \end{align*}
  In particular
  \begin{itemize}
    \item $\cH(A^0)$ is isomorphic to the vector space $\C^{(k-1)n(n-1)}$.
  \end{itemize}

  \begin{cor}
    The torus action of $\cH(A^0)$ changing the compatible framing
    corresponds to
    \begin{itemize}
      \item the conjugation action
        $t(\textbf{S})=(tS_1t^{-1},\dots,tS_{2k-2}t^{-1})$
        \begin{itemize}
          \item on the Stokes matrices
        \end{itemize}
    \end{itemize}
    and so, there is a bijection
    \[
      \Syst(A^0)/G\{z\}\cong(U_+\times U_-)^{k-1}/T
    \]
  \end{cor}
  \begin{cor}
    If $\Phi_0$ is contiuned once around $0$ in a positive sense, then on
    return to $\Sect_0$ is will become
    \[
      \Phi_0\cdot PS_{2k-2}\cdots S_2S_1P^{-1}M_0
    \]
    where
    \begin{itemize}
      \item $M_0=e^{2\pi\sqrt{-1}\cdot\Lambda^0}$ is the formal monodromy.
    \end{itemize}
  \end{cor}
\end{thm}


\begin{comment}
  \begin{rem}[1.41 from \cite{thboalch} on pages 16f]
    Note that in most of the recent references we have used, Stokes matrices
    are used to classify
    \begin{itemize}
      \item meromorphic connections within fixed formal \textbf{meromorphic
        classes, modulo meromorphic equivalence}.
    \end{itemize}
    Whereas here we classify
    \begin{itemize}
      \item meromorphic connections within fixed \textbf{formal analytic
        classes, modulo analytic equivalence},
    \end{itemize}
    as is done in the older literature.  The fact is that the sets equivalence
    classes are the same in both cases. It is important for us to work with
    analytic, rather than meromorphic gauge transformations, because then the
    $\C^\infty$ viewpoint in Chapter 3 is cleaner. This distinction relates to
    the difference between \textbf{‘regular singular’} connections and
    \textbf{‘logarithmic’} connections.
  \end{rem}
\end{comment}
%}}}

\section{Sheaf view: \cite{sabbah2007isomonodromic}} %{{{
Let
\begin{itemize}
  \item $X$ be the \textcolor{green!40!black}{parameter space}
    \begin{itemize}
      \item an analytic manifold equipped with coordinates $x_1,\ldots,x_n$.
    \end{itemize}
  \item $\sM^{good}$ be a meromorphic bundle
    \begin{itemize}
      \item on $D\times X$
      \item with poles along $\{0\}\times X$
    \end{itemize}
    equipped with \textbf{flat} connection
    \[
      \nabla^{good}:\sM\to\Omega_{D\times X}^1\otimes\sM
    \]
\end{itemize}
We will assume that $\left(\sM^{good},\nabla^{good}\right)$ is
\begin{itemize}
  \item a good model
    \begin{itemize}
      \item in the neighbourhood of any point $x^0\in X$,
    \end{itemize}
    that is that, for every $x^0$,
    \begin{itemize}
      \item there exist pairwise distinct germs $\phi_1,\dots,\phi_p\in
        t^{-1}\cO_{X,x^0}[t^{-1}]$ and
      \item nonzero germs of systems with regular singularity
        $\cR_{\phi_1},\dots,\cR_{\phi_p}$ along $\{0\}\times X$ such that
            we have, in the neighbourhood of $x^0$,
        \begin{itemize}
          \item an isomorphism $\sM_{x^0}^{good}
            \cong\bigoplus_k\left(\cE^{\phi_k}\otimes\cR_{\phi_k}\right)$.
            \begin{itemize}
              \item \TODO[wo/wie kommt hier die Umgebung rein? (p110)]
            \end{itemize}
        \end{itemize}
    \end{itemize}
\end{itemize}
\begin{comment}
  $k\neq l$ $\overset{\text{II.5.6}}{\Rightarrow{}}$ the order of the pole with
  respect to $t$ of $(\phi_k-\phi_l)(x,t)$ does not depend on $x$ in a
  neighbourhood of $x^0$.
\end{comment}

\subsection{The sheaf $\Aut^{<X}(\sM^{good})$}
\begin{defn}
  Define $\Aut^{<\{0\}\times X}(\tilde\sM^{good})$ or
  $\Aut^{<X}(\tilde\sM^{good})$ as the sheaf
  \begin{itemize}
    \item on $S^1\otimes X$
    \item of automorphisms of $\tilde\sM^{good}:=\sA_{\tilde D\times
      X}\otimes_{\cO_{D\times X}}\sM^{good}$
      witch are
      \begin{itemize}
        \item compatible with the connection and
        \item are formally equal to the identity\footnote{i.e., which
          induce the identity on $\hat \sM^{good}:=\hat \cO_{D\times
          X}\otimes_{\cO_{D\times X}}\sM^{good}$}.
      \end{itemize}
  \end{itemize}
  The sections of $\Aut^{<X}(\tilde\sM^{good})$ are called \emph{stokes
  matrices}.
\end{defn}
The germs $\phi_k$ of $\Aut^{<X}(\tilde\sM^{good})$ are polynomials in $t^{-1}$
without constant term, with coefficients in $\cO_{X,x^0}$.

An automorphism $a$ of $\tilde\sM^{good}$
\begin{itemize}
  \item on $I\times U$ where
    \begin{itemize}
      \item $I$ is an interval and
      \item $U$ is a sufficiently small neighbourhood of $x^0$
    \end{itemize}
\end{itemize}
can be decomposed into blocks
\[
  a_{kl}:\tilde\sE^{\phi_k}\otimes\tilde\sR_k
    \to
    \tilde\sE^{\phi_l}\otimes\tilde\sR_l
\]
The term thus takes the form $e^{\phi_k-\phi_l}b_{kl}$, where
\begin{itemize}
  \item $b_{kl}$ is a homomorphism of bundles with connection
    $\tilde\sR_k\to\tilde\sR_l$.
\end{itemize}
Assume that $X$ is a ball.
\begin{defn}
  Define
  \begin{itemize}
    \item $V_k$ to be the $\C$-vector space consisting of
      \begin{itemize}
        \item \textbf{multivalued horizontal sections}
          \begin{itemize}
            \item of $\tilde\sR_k$
            \item on $D^*\times X$
          \end{itemize}
      \end{itemize}
  \end{itemize}
  \begin{rem}
    \begin{itemize}
      \item $V_k$ is finite dimensional and
      \item on an open set like $]0,r[\times I\times X$,
        \begin{itemize}
          \item with $I\neq S^1$,
        \end{itemize}
        the \textbf{space of horizontal sections of $\bf\tilde\sR_k$}
        can be identified\footnote{This identification depending on the choice
        of a determination of the logarithm on $I$.} to \textbf{$\bf V_k$}.
    \end{itemize}
  \end{rem}
\end{defn}
Use the identification and $b_{kl}$ to induce a $\C$-linear map
\[
  b_{kl}:V_k\to V_l \,.
\]
\begin{rem}
  \begin{itemize}
    \item The matrix of $b_{kl}$
      \begin{itemize}
        \item with respect to the bases of $V_k$, $V_l$
      \end{itemize}
      is thus \textbf{constant}.
    \item The matrix of $b_{kl}$
      \begin{itemize}
        \item with respect to the $\cO_{D\times X}[t^{-1}]$-bases of $\sR_k$,
          $\sR_l$
      \end{itemize}
      has \textbf{moderate growth}.
  \end{itemize}
\end{rem}
For $k\neq l$, write
\[
  (\phi_k-\phi_l)(t,x)=\frac{\psi_{kl}(x)}{t^{n_{kl}}}\cdot U_{kl}(t,x)
\]
where
\begin{itemize}
  \item $n_{kl}>0$,
  \item $\psi_{kl}(x)\in\cO_{X,x^0}^*$ (i.e., does not vanish) and
  \item $u(t,x)$
    \begin{itemize}
      \item is holomorphic in the neighbourhood of $(0,x^0)$ and
      \item $u(0,0)=1$.
    \end{itemize}
\end{itemize}
We also choose some $C^\infty$ determination $\eta_{kl}(x)$ of the argument of
$\psi_{kl}(x)$:
\[
  \psi_{kl}(x)=|\psi_{kl}(x)|\cdot e^{i\eta_{kl}(x)}
\]
\begin{rem}
  $\psi_{kl}(x)=-\psi_{lk}(x)$
\end{rem}
\begin{lem}
  Let us pick $e^{i\theta^0}\in S^1$
  \begin{enumerate}
    \item For $k\neq l$, the matrix $e^{\phi_k-\phi_l}b_{kl}$
      \begin{itemize}
        \item in $\cO_{D\times X}[t^{-1}]$-bases of $\sE^{\phi_k}\otimes\sR_k$,
          $\sE^{\phi_l}\otimes\sR_l$
      \end{itemize}
      has enties in
      \begin{itemize}
        \item $\sA_{S^1\times X}^{<X}$
          \begin{itemize}
            \item in the neighbourhood of $(e^{i\theta^0},x^0)$
          \end{itemize}
      \end{itemize}
      if and only if $b_{kl}=0$ or
      $\cos(n_{kl}\theta^0-\eta_{kl})<0$.
    \item The matrix $b_{kk}-\Id$ has entries in
      \begin{itemize}
        \item $\sA_{S^1\times X}^{<X}$
          \begin{itemize}
            \item in the neighbourhood of $(e^{i\theta^0},x^0)$
          \end{itemize}
      \end{itemize}
      if and only if $b_{kk}-\Id=0$.
  \end{enumerate}
  \begin{comment}
    \begin{proof}[Proofsketch]
      For the first point, it is enough to check that $e^{\phi_k-\phi_l}$
      belongs to $\sA_{S^1\times X}^{<X}$ in the neighbourhood of
      $(e^{i\theta^0},x^0)$ if and only if $\Re(\phi_k-\phi_l)<0$ on a
      sufficiently small neighbourhood of $(e^{i\theta^0},x^0)$, then to
      express this condition on the leading term of $\phi_k-\phi_l$.  Then, one
      notices that, as $b_{kl}$ has moderate growth, it does not affect the
      rapid decay or the exponential growth property.
    \end{proof}
  \end{comment}
\end{lem}
Let
\begin{itemize}
  \item $V=\bigoplus_kV_k$ be the $\C$-vector space of multivalued horizontal
    sections of $\bigoplus_k\sR_k$ on $D^*\times X$.
  \item $\sL$ the subsheaf of the constant sheaf $\Aut(V)$
    \begin{itemize}
      \item of which the germ at $(e^{i\theta^0},x^0)$ is
        the space of automorphism $\Id+(\oplus c_{kl})$ with
        \begin{itemize}
          \item $c_{kk}=0$ and
          \item $c_{kl}:V_k\to V_l$ is nonzero if and only if
            $\cos(n_{kl}\theta^0-\eta_{kl})<0$.
        \end{itemize}
    \end{itemize}
\end{itemize}
\begin{rem}
  Any Stokes matrix (local section of $\sL$) is
  \begin{itemize}
    \item unipotent
      \begin{itemize}
        \item if we suitably order the set of indices, we can assume that
          $c_{kl}=0$ for $k<l$ and the matrix of $\oplus c_{kl}$ is strictly
          uppertriangular.
      \end{itemize}
  \end{itemize}
\end{rem}
\begin{cor}[The Stokes martices are constant]
  For
  \begin{itemize}
    \item any open interval $I\neq S^1$,
  \end{itemize}
  the restriction to $I\times X$ of the
  sheaves $\bf\Aut^{<X}(\tilde\sM^{good})$ and $\bf\sL$ are isomorphic:
  \[
    \Aut^{<X}(\tilde\sM^{good})_{|I\times X}
    \cong
    \sL_{|I\times X}
  \]
\end{cor}
\begin{cor}[Base change for horizontal sections]
  Let
  \begin{itemize}
    \item $i:S^1\times\{x^0\}\hookrightarrow S^1\times X$ denote the inclusion
  \end{itemize}
  then the natural morphism
  \[
    i^{-1}\Aut^{<X}(\tilde\sM^{good})\to\Aut^{<0}(\widetilde{i^+\sM}^{good})
  \]
  is an isomorphism.
\end{cor}
\begin{defn}
  The \emph{manifold of Stokes directions}
  \begin{itemize}
    \item in $S^1\times U$, where
      \begin{itemize}
        \item $U$ is a neighbourhood of $x^0$ on which the $\phi_k$ are defined
          and satisfy the goodness properties
      \end{itemize}
  \end{itemize}
  is the union of the sets having the equation
  \[
    \cos(n_{kl}\theta-\eta_{kl}(x))=0 \,.
  \]
  The restriction of this set to $S^1\times\{x^0\}$ is a finite set of points,
  called the \emph{Stokes directions at $x^0$}.
  \begin{rem}
    A pair $(k,l)$ contributes to the directions
    \[
      \frac{\left(\eta_{kl}+\frac{\pi}{2}+j\pi\right)}{n_{kl}}
    \]
    for $j=0,\dots,2n_{kl}-1$.
  \end{rem}
\end{defn}
\begin{ex}[Straightening the Stokes manifold (\cite{sabbah2007isomonodromic}
  6.8)]
\end{ex}

\subsection{The Stokes sheaf $\St_X(\sM^{good})$}
\begin{defn}
  Define the \emph{stokes sheaf} $\St_X(\sM^{good})$ as the sheaf (on $X$)
  associated to the presheaf
  \[
    U\mapsto H^1(S^1\times U,\Aut^{<X}(\tilde \sM^{good}))\,.
  \]
\end{defn}
\begin{thm}[The stokes sheaf is locally constant (6.1)]
  The Stokes sheaf $\St_X(\sM^{good})$ is
  \begin{itemize}
    \item a locally constant
  \end{itemize}
  sheaf of pointed sets.
  \begin{comment}
    \begin{proof}
      \cite[p. 113]{sabbah2007isomonodromic}: paragraph 6.c.
    \end{proof}
  \end{comment}
\end{thm}
\begin{prop}[Base change for the Stokes sheaf]
  The germ at $x^0$ of the Stokes sheaf $\St_X(\sM^{good})$ is equal to the
  Stokes space\footnote{Is equipped with a natural structure of affine space.}
  $H^1(S^1,\Aut^{<0}(\widetilde{i^+\sM}^{good}))$ of the connection
  $i^+\sM^{good}$.
\end{prop}

\subsection{The sheaf $\sH_X$}
\begin{defn}
  Let $(\sM,\nabla)$ be
  \begin{itemize}
    \item a meromorphic bundle
      with
      \begin{itemize}
        \item flat connection,
          \begin{itemize}
            \item having poles along $\{0\}\times X$.
          \end{itemize}
      \end{itemize}
  \end{itemize}
  \begin{defn}
    $(\sM^{good},\nabla^{good})$ is a \emph{good formal model} for
    $(\sM,\nabla)$ if
    \begin{itemize}
      \item there exists an isomorphism
        \[
          \hat f:(\widehat\sM,\widehat\nabla)\overset{\sim}{\to}
          (\widehat{\sM^{good}},\widehat{\nabla^{good}})
        \]
        \begin{itemize}
          \item of sheafes of $\hat\cO_{D\times X}$-modules
          \item compatible with connection.
        \end{itemize}
    \end{itemize}
  \end{defn}
  \begin{defn}
    Two germs $(\sM,\nabla,\hat f)$ and $(\sM',\nabla',\hat f')$ along
    $\{0\}\times X$ are \emph{isomorphic}, if
    \begin{itemize}
      \item there exists an isomorphism\footnote{such an isomorphism then is
        unique} $g:(\sM,\nabla)\to(\sM',\nabla')$ with
        \begin{itemize}
          \item $\hat f=\hat f'\circ\hat g$.
        \end{itemize}
    \end{itemize}
  \end{defn}
  Let
  \begin{itemize}
    \item $(\sM^{good},\nabla^{good})$ be a fixed good formal model.
  \end{itemize}
  Define
  \begin{itemize}
    \item $\sH_X$ be the (pre-) sheaf\footnote{Lemma 6.2: it is a sheaf.} where
      \begin{itemize}
        \item $\sH_X(U)$ is the (pointed) set
          \begin{itemize}
            \item of isomorphism classes of
              germs\footnote{along $\{0\}\times U$}
              $(\sM,\nabla,\hat f)$ defined on $U$,
              \begin{comment}
                \begin{itemize}
                  \item Why germs??
                  \item along $\{0\}\times U$??
                \end{itemize}
              \end{comment}
            \item equipped with a distinguished element, namely, the class of
              ${(\sM^{good},\nabla^{good},\widehat{\Id})}_{|U}$.
          \end{itemize}
      \end{itemize}
  \end{itemize}
\end{defn}

\subsection{Main Result}
Use Thoreme~\ref{thm:sectDecomp} to construct\footnote{see:
\cite[p. 111f]{sabbah2007isomonodromic}} a homomorphism of sheaves of
pointed sets
\[
  \sH_X\to\St_X(\sM^{good})
\]
\begin{itemize}
  \item If \begin{itemize}
      \item $(\sM,\nabla,\hat f)$ is defined on $U$
    \end{itemize}
    there exists
    \begin{itemize}
      \item an open covering $\mathfrak{W}$ of $S^1\times U$ and
      \item for any $W_i$ in $\mathfrak{W}$
        \begin{itemize}
          \item an isomorphism
            $f_i:(\tilde\sM,\tilde\nabla)_{W_i}
              \overset{\sim}{\to}
              (\tilde\sM^{good},\tilde\nabla^{good})_{W_i}$
        \end{itemize}
    \end{itemize}
    such that $\hat f_i=\hat f$.

    Then $(f_jf_i^{-1})_{i,j}$ is a cocycle of the sheaf
    $\Aut^{<X}(\tilde\sM^{good})$ relative to the covering $\mathfrak{W}$.
    \begin{itemize}
      \item If $f_i'$ is another lifting of $\hat f$\dots
      \item If $(\sM,\nabla,\hat f)$ and $(\sM',\nabla',\hat f')$ are
        isomorphic\dots
    \end{itemize}
    We have thus defined a mapping of pointed sets
    \[
      \Gamma(U,\sH_X)\to H^1(S^1\times U,\Aut^{<X}(\tilde\sM^{good}))
    \]
    from which we deduce the homomorphism of sheaves.
\end{itemize}

\begin{thm}[The Stokes sheaf classifies meromorphic connections with fixed
  normal type (6.3)]
  The homomorphism so defined
  \[
    \sH_X\to St_X(\sM^{good})
  \]
  is an isomorphism of sheaves of pointed sets.
\end{thm}
\begin{cor}[Analytic extension with fixed formal structure]
  Let
  \begin{itemize}
    \item $X$ be $1$-connected and
    \item $(\sM^0,\nabla^0,\hat f^0)$ be a germ of meromorphic bundle
      \begin{itemize}
        \item on $D$
        \item with pole at $0$
        \item equipped with a formal isomorphism
          \[
            \hat f^0:(\sM^0,\nabla^0)
              \overset{\sim}{\to}
              i^+(\sM^{good},\nabla^{good})
          \]
          where
          \begin{itemize}
            \item $i$ denotes the inclusion
              $D\times\{x^0\}\hookrightarrow D\times X$
          \end{itemize}
      \end{itemize}
  \end{itemize}
  then there exists
  \begin{itemize}
    \item a meromorphic bundle with connection $(\sM,\nabla,\hat f)$
      \begin{itemize}
        \item equipped with a formal isomorphism $\hat f$ to
          $(\sM^{good},\nabla^{good})$
      \end{itemize}
  \end{itemize}
  such that
  \begin{itemize}
    \item $i^+(\sM,\nabla,\hat f)$ is isomorphic to
      $(\sM^0,\nabla^0,\hat f^0)$.
  \end{itemize}
  Such an object is unique up to unique isomorphism.
\end{cor}
%}}}

\section{Comparison}
\subsection{Situation} %{{{
\begin{paracol}{2} %%%%%%%%%%%%%%%%%%%%%%%%%%%%%%%%%%%%%%%%%%%%%%%%%%%%%%%%%%%%
  Fix the data $\textbf{a}$ of
  \begin{itemize}
    \item a Divisor $D=\sum k_i(a_i)$ on $\P^1$ and
    \item and connection germs $d-{}^iA^0$ at each $a_i$.
  \end{itemize}
  Let $(V,\nabla,\textbf{g})$ be a compatibly framed meromorphic connection on
  a holomorphic vector bundle $V\to\P^1$ with a irregular type $\textbf{a}$.

  \switchcolumn{} %%%%%%%%%%%%%%%%%%%%%%%%%%%%%%%%%%%%%%%%%%%%%%%%%%%%%%%%%%%%%
  Let
  \begin{itemize}
    \item $X$ be the parameter space
      \begin{itemize}
        \item an analytic manifold equipped with coordinates $x_1,\ldots,x_n$.
      \end{itemize}
    \item $\sM$ be a meromorphic bundle
      \begin{itemize}
        \item on $D\times X$
        \item with poles along $\{0\}\times X$
      \end{itemize}
      with \textbf{flat} connection
      \[
        \nabla:\sM\to\Omega_{D\times X}^1\otimes\sM
      \]
    \item $\cM$ denote its germ at $(0,x^0)\in D\times X$
      \begin{itemize}
        \item this is a module over the ring
          \[
            \cO_{D\times X,(0,x^0)}[t^{-1}]=\C\{t,x_1,\dots,x_n\}[t^{-1}]
          \]
      \end{itemize}
  \end{itemize}
\end{paracol} %%%%%%%%%%%%%%%%%%%%%%%%%%%%%%%%%%%%%%%%%%%%%%%%%%%%%%%%%%%%%%%%%
%}}}

